\documentclass[10pt,a4paper]{amsart}
\usepackage[margin=19mm]{geometry}
\usepackage{amsmath,amssymb,mathtools}
\usepackage[scr=rsfs,cal=euler]{mathalfa}
\usepackage{xfrac}
\usepackage[unicode,hidelinks,bookmarks=false]{hyperref}
\newcommand{\ie}{i.e.\ }
\newcommand{\cf}{cf.\ }
\newcommand{\resp}{respectively\ }
\newcommand{\Subsection}{\subsection}
\providecommand{\nobreakdash}{\nobreak\hbox{-}}
\setlength{\emergencystretch}{2em}
\multlinegap0pt

\theoremstyle{plain}
\newtheorem{theorem}{Theorem}[section]
\newtheorem{lemma}[theorem]{Lemma}
\newtheorem{proposition}[theorem]{Proposition}
\newtheorem{corollary}[theorem]{Corollary}
\newtheorem{conjecture}[theorem]{Conjecture}
\newtheorem{definition}[theorem]{Definition}
\newtheorem{example}[theorem]{Example}
\newtheorem{remark}[theorem]{Remark}
\newtheorem{question}[theorem]{Open Problem}
\newtheorem{problem}[theorem]{Open Problem}
\title[Cyclic Latin Eulerian Numbers]{Cyclic Latin Eulerian Numbers}
\author{Madjid Mirzavaziri, Daniel Yaqubi, Daniel Yaqubi$^{*}$}
\date{Source: arXiv:2609.28808v1 (CC BY 4.0)}
\begin{document}
\maketitle

\noindent\emph{Source and licence.} The delivered work is the article shown in the title above, version arXiv:2609.28808v1 (CC BY 4.0), distributed under the Creative Commons Attribution 4.0 International licence, as recorded in the arXiv licence metadata for this exact version and shown on the abstract page saved with this item. The full licence notice, which carries the licence URL and the address of that abstract page, is the bundled file \texttt{licenses/arxiv-2609.28808-CC-BY-4.0.txt}.
\par\smallskip

\noindent\textit{Editorial scope.} Counted unit: the article's Theorem (source label thm:cyclic-boundary-ascents) (source0.tex lines 228--238, source label \texttt{thm:cyclic-boundary-ascents}): ``For every , the lower boundary values of are T n c(n-1) n, T n c(n) 0, and T n c(n 1) n. By reversal symmetry, the upper boundary values satisfy T n c ((n-1) 2 ) n, T n c ((n-1) 2-1 ) 0, and T n c ((n''. It is delivered together with the article's complete original proof (source0.tex lines 240--252, one proof environment adjacent to the statement, 13 lines), copied line for line from the article's own text. Required context retained uncounted: prop:cyclic-symmetry-divisibility (lines 214--220). Editorial formatting only: an AMS \texttt{amsart} preamble that restates the article's own macro definitions and its own theorem declarations, and the theorem environments the retained lines use; the article's own class file is neither shipped nor input; the retained environments are renumbered sequentially in order of appearance, whereas the article prints them with its own numbering; the article's equation numbering is not reproduced and every cross-reference inside the retained text resolves to the wrapper's numbering. No citation occurs inside any retained line, so the wrapper carries no bibliography; All retained bodies are exact copies of the bundled \texttt{sources/211620/source0.tex}, so no omission receipt applies. No asset-inclusion command occurs inside any retained span and the package ships no image asset.


\section*{Retained context: prop:cyclic-symmetry-divisibility (source0.tex lines 214--220)}

\begin{proposition}
\label{prop:cyclic-symmetry-divisibility}
For all $n \ge 1$ and $m \in \mathbb{Z}_{\ge 0}$,
\begin{equation}
T_n^c(m) = T_n^c\left(n(n-1) - m\right) \qquad \text{and} \qquad n \mid T_n^c(m).
\end{equation}
\end{proposition}


\section{The counted unit: Theorem (source label thm:cyclic-boundary-ascents) and its complete original proof}

\begin{theorem}
\label{thm:cyclic-boundary-ascents}
For every $n \ge 3$, the lower boundary values of $T_n^c(m)$ are
\begin{equation}
T_n^c(n-1) = n, \qquad T_n^c(n) = 0, \qquad \text{and} \qquad T_n^c(n+1) = n.
\end{equation}
By reversal symmetry, the upper boundary values satisfy
\begin{equation}
T_n^c\left((n-1)^2\right) = n, \qquad T_n^c\left((n-1)^2-1\right) = 0, \qquad \text{and} \qquad T_n^c\left((n-1)^2-2\right) = n.
\end{equation}
\end{theorem}

\begin{proof}
Setting $e_j := n - 1 - \delta_j \ge 0$, the total ascent statistic can be written as $\Sigma(L_\pi) = (n-1) + \sum_{j=1}^{n-1} e_j$.

\begin{itemize}
    \item For $\Sigma = n-1$, we must have $e_j = 0$ ($\delta_j = n-1$) for all $j$. The unique difference sequence $(n-1, \dots, n-1)$ defines a cyclic descending progression with exactly $n$ translations, so $T_n^c(n-1) = n$.
    
    \item For $\Sigma = n$, we have $\sum e_j = 1$, meaning exactly one difference is $n-2$ while the remaining $n-2$ differences are $n-1$. The sum of all differences then satisfies $\sum_{j=1}^{n-1} \delta_j = (n-2)(n-1) + (n-2) \equiv 0 \pmod n$, forcing $\pi_n = \pi_1$, which contradicts $\pi \in S_n$. Thus $T_n^c(n) = 0$.
    
    \item For $\Sigma = n+1$, we have $\sum e_j = 2$. If $e_j = 2$ for some $j$ ($\delta_j = n-3$), partial sum evaluation again forces a repeated residue in $\pi$. Thus $e_r = e_s = 1$ ($r < s$), so $\delta_r = \delta_s = n-2$. To prevent repeated entries in $\pi$, the indices are uniquely forced to $r = 1$ and $s = n-1$. The resulting unique difference pattern $(n-2, n-1, \dots, n-1, n-2)$ admits $n$ translations, yielding $T_n^c(n+1) = n$.
\end{itemize}

Applying the reversal symmetry of Proposition~\ref{prop:cyclic-symmetry-divisibility} directly completes the upper boundary evaluations.
\end{proof}

\end{document}
